\documentclass[11pt,a4paper]{article}
\usepackage[margin=25mm]{geometry}
\usepackage{amsmath,amssymb,amsthm,hyperref}
\setlength{\emergencystretch}{2em}
\newtheorem{theorem}{Theorem}
\newtheorem{corollary}[theorem]{Corollary}
\newcommand{\gen}[1]{\langle#1\rangle}
\title{Parity of a paired marking with two endpoint anchors}
\author{Human Compression project}\date{}
\begin{document}\maketitle
We acquire the separately proved anchor-deletion theorem and the exact
seven- and nine-point boundaries, including their distinct reviews.
The infinite deduction below uses these actual action statements;
no general near-half classification is taken as a premise.

Let $G=\gen{h_0,\ldots,h_{r-1}}$ act transitively on $2r+1$ points.
The marks form a string C-marking. Every colour is paired: it has
at least two edges with endpoints in different orbits of the
corresponding maximal omission. A first anchor $b$ is moved by
$h_0$ and fixed by all other marks. A last anchor $c$ is defined
in the reversed order, and $b,c$ are distinct.

The acquired anchor-deletion theorem says that for $r\ge5$ the
actual marks begin
\[
 h_0=(b\ x)(u\ v),\qquad h_1=(x\ u)(v\ z)\beta,       \tag{1}
\]
where $\beta$ is supported in $V\setminus\{z\}$ and
$V=\Omega\setminus\{b,x,u,v\}$. The actual tail
$K=\gen{h_2,\ldots,h_{r-1}}$ fixes the four removed points,
is transitive on $V$, and retains the paired property, the
intersection property and the two anchors $z,c$. Its degree and
rank are $2(r-2)+1$ and $r-2$. For $r\ge7$, the same theorem
proves $\beta=1$ by an actual subgroup-intersection argument.
These conclusions do not assume that $G$ or $K$ is symmetric.

\begin{theorem}\label{even}
Every mark of $G$ is even. In particular, no natural symmetric
marking of degree $2r+1$ and rank $r$ can have the stated paired
property and two endpoint anchors.
\end{theorem}
\begin{proof}
We use induction on $r$. Rank one is impossible, since a single
involution cannot be transitive on three points. At rank two each
mark on five points has at least two transpositions, by eligibility,
and at most two by the matching property. Thus both are even.
At rank three the acquired seven-point anchor theorem excludes
transitivity already from the string and anchor conditions.

The nine-point boundary supplies the complete rank-four result.
Every transitive paired string marking with the two anchors is,
after a point relabelling respecting the anchor order, the path
whose successive edge labels are
\[
                         0,1,0,1,2,3,2,3.               \tag{2}
\]
In particular each mark is a double transposition and the generated
group is $A_9$. Its intersection property and actual group are
independently checked in that acquired boundary, although its
classification of the transitive paired tuples does not require IP.

At rank five, (1) would give a rank-three seven-point tail with
two anchors, which is impossible. Consider rank six. Its tail has
rank four and degree nine, so the acquired literal path (2) applies.
Use its point labels so the first tail mark is
\[
 h_2=(z\ p)(q\ w).
\]
The tail omission $D=\gen{h_3,h_4,h_5}$ has orbits
$\{z\},\{p,q\}$ and a six-point orbit containing $c$.
Its centralizer fixing $z,c$ is exactly $\{1,(p\ q)\}$.
The factor $\beta$ in (1) belongs to this centralizer: it fixes
$z,c$, and $h_1$ commutes with $h_3,h_4,h_5$, which fix
$x,u,v,z$. Thus $\beta$ is either $1$ or $(p\ q)$.

In the latter case $(h_0h_1)^5=\beta$, since the other factor
of this product is a five-cycle on $b,x,u,v,z$. Conjugation
by $h_2$ gives $(z\ w)$. Therefore
\[
 \delta=(p\ q)(z\ w)\in\gen{h_0,h_1,h_2}\cap K.
\]
The second membership uses the actual $K=A_9$, since $\delta$
is even and supported in its nine-point orbit. The original IP
makes this intersection $\gen{h_2}$, whereas $\delta$ is
neither its identity nor its nonidentity element. This contradiction
gives $\beta=1$. Both marks in (1) and every tail mark are now even.

For $r\ge7$, the acquired theorem already gives $\beta=1$.
Induction on the actual tail shows that all its marks are even.
The two remaining marks in (1) are double transpositions, proving
the assertion. Both inductions decrease rank by two on the exact
support; no projected subgroup is assigned an intersection property.
\end{proof}

\begin{corollary}\label{critical}
Under the strict smaller-degree sufficient-split induction, the
critical overlap of a six-point rank-three even $A_5$ factor and
a fractured natural symmetric factor of even degree $m$ and rank
$m/2+1$ has an internal sufficient split.
\end{corollary}
\begin{proof}
We use the exact configuration and definition of sufficient split
in the acquired first-cut and endpoint manuscripts. The full
degree is $n=m+5$. The first anchored unique cut in the symmetric
factor is perfect. If its index is positive, the complete first-cut
theorem supplies a sufficient split, using only the stated strict
smaller-degree induction. If its index is zero, the complete
endpoint theorem either supplies such a split or produces an actual
natural symmetric middle of degree $2k+1$ and rank $k$, paired
in every colour and with two distinct endpoint anchors. Theorem
\ref{even} excludes precisely that residual. Thus both routes
finish. The rank induction of Theorem~\ref{even} is independent
of the outer sufficient-split induction, so no circular premise
is introduced.
\end{proof}

\paragraph{Attribution and scope.}
The source is Cameron--Fernandes--Leemans, \emph{The number of
string C-groups of high rank}, arXiv:2212.12723v2. Its paired
near-half classification includes (2). Here only the exact
nine-point boundary is enumerated; the infinite parity deduction
uses the separately proved anchor deletion and actual tail.
The classical rank results and earlier support arguments remain
attributed through their acquired proofs. This component excludes
the indicated middle and closes a conditional sufficient-split
route. The complete degree induction, the stable counting
correspondence, contraction and reversal compatibility remain
separate obligations. No whole-original, conjecture, Lean or
compression credit is asserted.
\end{document}
